\clearpage\nobalance
\section*{Author Notes}
\subsection*{Acknowledgements}
I thank Prof. Luyao Zhang, Program Chair, Prof. Ken Rogerson, Program Discussant, Jiayang [SURNAME], Xuantong [SURNAME], and the COMSCI/ECON 206 Autumn 2026 Session 1 class. Jiayang's reputation discussion motivated the behavioral demand condition; Xuantong's rule-comparison discussion motivated completed access as an outcome. Prof. Zhang's feedback led me to compare fixed and chosen intermediary margins and include coordination costs. I am responsible for the proposal, code, and remaining errors.
\subsection*{Individual contribution}
I designed the farmer-intermediary application, sequential model, synthetic baseline, computational checks, and proposed reputation test. I wrote and verified the code and adapted the cited theories and class feedback with acknowledgement.
\bibliographystyle{ACM-Reference-Format}\bibliography{references}
\appendix
\section{Technical details and reproducibility}
The baseline uses $q(m,e)=12e\max\{0,4-m\}$, $g(e)=30e^2$, $w=6$, $c=4$, and $K=12$. Fixed $m=1$; chosen searches $m\in[0,4]$ by $0.05$ and $e\in[0,1]$ by $0.01$. The committed notebook and `companion/outputs/` reproduce the values in Section 3. All results are synthetic.
\subsection{AI-use disclosure}
OpenAI GPT-6 assisted with drafting, debugging, documentation, and consistency checks on September 13, 2026 after I selected the question and supplied course feedback. I verified equations, outputs, citations, and scope. Initial reasoning, handwritten reflection, and peer reviews are Human-Only.
\section{Cumulative Intellectual Development}
This proposal develops Intellectual Statements I and II through the \emph{2056 Nobel and Turing Laureate Program Launch Workshop}, September 7, 2026, Duke Kunshan University, IB 2050. My session was Group 1 and my role was economist. My first statement asked whether price information could strengthen farmers' bargaining positions. Week 2 shifted the mechanism toward pooled selling and organization. Prof. Zhang then specified one chain from selling rule to household opportunity: compare fixed and chosen margins and include coordination costs. I retained the bargaining concern but replaced broad transparency and inequality claims with measurable farmer earnings and completed household purchases.
\section{Field and open-source inspiration}
The \emph{Joint Interdisciplinary Field Trip---Technology, Computation, Culture \& Society in Action} occurred September 4, 2026. I attended the Tencent Shanghai Office and Shanghai Science and Technology Museum visits. Retain the two existing photographs with exact date and provenance in the figure captions; distinguish observation from interpretation. Open Food Network motivates coordinated producer listings~\cite{openfoodnetwork}; CKAN motivates publishing parameters and outputs~\cite{ckan}. These sources inspire the design but are not causal evidence for farmer earnings.
\section{Review response and revision record}
Peer review, rebuttal, reviewer follow-up, and v2 will be completed in the Human-Only class process. Preserve v1, both reviews, the response, follow-up, and v2 here; do not place private review documents in the public GitHub repository.
\clearpage\onecolumn
\section{Structured abstract and cumulative proposal map}\label{app:abstract}
\begin{table}[H]\caption{Appendix E: eight-row structured abstract.}\renewcommand{\arraystretch}{1.15}\begin{tabularx}{\textwidth}{@{}p{.25\textwidth}Y@{}}\toprule\textbf{Element}&\textbf{Current statement}\\\midrule
Background&Small farmers can pool products and use intermediaries to reach underserved households; margin and selection rules affect effort, prices, and surplus.\\
Motivation and application&A platform or cooperative must choose how an intermediary earns a margin without reducing farmer participation or completed household purchases.\\
Three connected questions&Economics compares rules; computation finds thresholds; behavioral science estimates how reputation changes trust and demand.\\
Method and model assumptions&Sequential game with farmer coordination cost, intermediary margin and effort, household purchase, grid search, sensitivity analysis, and a reputation experiment.\\
Results or expected results&Current synthetic baseline favors fixed margins for farmer earnings and completed units, while chosen margins favor intermediary profit; reversals remain expected under other parameters.\\
Intellectual merits&The project identifies the mechanism linking margin governance, coordination cost, computational thresholds, and trust.\\
Practical impacts&Platforms can compare rules before implementation; the Margin Rule Lab will teach the tradeoff for SDG 4.\\
Feedback received and revision made&Prof. Zhang recommended fixed versus chosen margins, coordination costs, and one chain to household opportunity. I replaced broad inequality claims with this bounded mechanism.\\\bottomrule\end{tabularx}\end{table}
